%%%%%%%%%%%%%%%%%%%%%%%%%%%%%%%%%%%%%%%%%%%%%%%%%%%%%%%%%%%%%%%%%%%%
%% RL for LLMs: A Systematic GRPO Scaling Study
%% Capstone Project — Group 6
%% MTech Data Science & AI, PES University, Bengaluru
%%
%% Chapter 2: Foundation (Literature Review)
%%%%%%%%%%%%%%%%%%%%%%%%%%%%%%%%%%%%%%%%%%%%%%%%%%%%%%%%%%%%%%%%%%%%

\documentclass[12pt,a4paper]{report}

%% ---------- Page geometry ----------
\usepackage[
  top=2.5cm,
  bottom=2.5cm,
  left=3cm,
  right=2.5cm
]{geometry}

%% ---------- Mathematics ----------
\usepackage{amsmath}
\usepackage{amssymb}

%% ---------- Typography & spacing ----------
\usepackage[utf8]{inputenc}
\usepackage[T1]{fontenc}
\usepackage{setspace}
\onehalfspacing

%% ---------- Graphics & figures ----------
\usepackage{graphicx}
\usepackage{tikz}
\usepackage{caption}
\usepackage{subcaption}

%% ---------- Tables ----------
\usepackage{booktabs}

%% ---------- Algorithms ----------
\usepackage[ruled,vlined,linesnumbered]{algorithm2e}

%% ---------- Code listings ----------
\usepackage{listings}
\lstset{
  basicstyle=\ttfamily\small,
  breaklines=true,
  frame=single,
  numbers=left,
  numberstyle=\tiny,
  captionpos=b
}

%% ---------- Hyperlinks ----------
\usepackage[
  colorlinks=true,
  linkcolor=blue!70!black,
  citecolor=green!50!black,
  urlcolor=blue!60!black
]{hyperref}

%% ---------- Bibliography ----------
\usepackage[round,sort&compress]{natbib}
\bibliographystyle{plainnat}

%% ---------- Title information ----------
\title{%
  \vspace{-1cm}
  \textbf{RL for LLMs: A Systematic GRPO Scaling Study} \\[1.5em]
  \large Capstone Project Report \\[0.8em]
  \large MTech in Data Science \& Artificial Intelligence
}
\author{%
  \textbf{Group 6} \\[1em]
  Department of Computer Science and Engineering \\
  PES University, Bengaluru --- 560085, India
}
\date{Academic Year 2025--2026}

%%%%%%%%%%%%%%%%%%%%%%%%%%%%%%%%%%%%%%%%%%%%%%%%%%%%%%%%%%%%%%%%%%%%
\begin{document}

%% ================================================================
%% TITLE PAGE
%% ================================================================
\begin{titlepage}
  \centering
  \vspace*{1.5cm}

  {\LARGE\bfseries PES University, Bengaluru}\\[0.4cm]
  {\large Department of Computer Science and Engineering}\\[0.3cm]
  {\large MTech in Data Science \& Artificial Intelligence}\\[2.5cm]

  \rule{\textwidth}{1.2pt}\\[0.5cm]
  {\Huge\bfseries RL for LLMs:\\[0.3em]
   A Systematic GRPO Scaling Study}\\[0.5cm]
  \rule{\textwidth}{1.2pt}\\[2cm]

  {\Large\textbf{Capstone Project Report}}\\[1.5cm]

  {\large\textbf{Group 6}}\\[2.5cm]

  \vfill
  {\large Academic Year 2025--2026}
\end{titlepage}

%% ================================================================
%% CHAPTER 2: FOUNDATION
%% ================================================================
\setcounter{chapter}{1}          % so that \chapter gives "Chapter 2"
\chapter{Foundation}
\label{ch:foundation}

%% ----------------------------------------------------------------
\section{Introduction}
\label{sec:foundation-intro}

Large language models (LLMs) have demonstrated remarkable capabilities across
a broad spectrum of natural-language tasks, yet their raw pre-trained behaviour
often diverges from the nuanced requirements of downstream applications.
Bridging this gap---aligning model outputs with human intent, factual
correctness, and structured reasoning---has become one of the central
challenges of modern NLP research.  The present chapter lays the theoretical
and empirical groundwork upon which our scaling study is built, surveying the
key ideas, algorithms, and benchmarks that collectively motivate and inform
our experimental design.

\paragraph{Reinforcement Learning from Human Feedback (RLHF).}
The RLHF pipeline, popularised by InstructGPT \citep{ouyang2022instructgpt},
established a three-stage recipe---supervised fine-tuning, reward-model
training, and reinforcement-learning-based policy optimisation---that
transformed how practitioners align LLMs with human preferences.  We begin
in Section~\ref{sec:rlhf-pipeline} by reviewing this pipeline in detail,
since every subsequent algorithmic advance discussed in this chapter can be
understood as an attempt to simplify, stabilise, or scale one or more of
its stages.

\paragraph{Policy optimisation: from PPO to DPO to GRPO.}
Proximal Policy Optimisation \citep[PPO;][]{schulman2017ppo} served as the
workhorse optimiser within early RLHF systems, but its reliance on a
separately trained value network and careful hyperparameter tuning introduced
significant computational overhead.  Direct Preference Optimisation
\citep[DPO;][]{rafailov2023dpo} later showed that the reward-modelling
step could be bypassed entirely by re-parameterising the Bradley--Terry
preference objective in terms of the policy itself.  Most recently, Group
Relative Policy Optimisation \citep[GRPO;][]{shao2024deepseekmath} proposed
an elegant middle ground: it retains the explicit reward signal of PPO but
eliminates the value function by computing advantages relative to a group
of sampled completions, thereby reducing memory requirements and simplifying
the training loop.  The DeepSeek-R1 effort \citep{deepseek2025r1}
subsequently demonstrated that GRPO can elicit sophisticated multi-step
reasoning in large-scale models.  Section~\ref{sec:policy-optimisation}
traces this algorithmic lineage, formalises each objective, and analyses the
trade-offs that make GRPO an attractive choice for resource-constrained
scaling experiments.

\paragraph{Parameter-efficient fine-tuning.}
Full fine-tuning of models with billions of parameters is prohibitively
expensive for most research groups.  Low-Rank Adaptation
\citep[LoRA;][]{hu2021lora} offers a compelling alternative by injecting
trainable low-rank matrices into frozen pre-trained weights, dramatically
reducing the number of trainable parameters while preserving expressive
power.  In Section~\ref{sec:peft}, we review LoRA and its variants,
discussing how they interact with reinforcement learning objectives and
why they are essential enablers of the scaling experiments presented in
later chapters.

\paragraph{Mathematical reasoning benchmarks.}
Rigorous evaluation requires benchmarks with unambiguous ground-truth
answers.  We focus on two widely adopted mathematical reasoning datasets:
GSM8K \citep{cobbe2021gsm8k}, a collection of grade-school math word
problems, and MATH \citep{hendrycks2021math}, a more challenging suite
spanning algebra, number theory, geometry, and beyond.  Together with
chain-of-thought prompting strategies \citep{wei2022chainofthought}, these
benchmarks provide a controlled test-bed for measuring whether
reinforcement-learning-based fine-tuning genuinely improves a model's
capacity for step-by-step reasoning rather than superficial pattern
matching.  Section~\ref{sec:benchmarks} describes these datasets, their
evaluation protocols, and the metrics we adopt.

\paragraph{Related scaling studies.}
Finally, Section~\ref{sec:related-scaling} surveys prior work that has
investigated how RLHF and its variants behave as a function of model
size, dataset scale, and compute budget.  These studies provide important
baselines and methodological precedents for our own systematic exploration
of GRPO scaling behaviour.

\medskip
\noindent
Taken together, the topics covered in this chapter establish the conceptual
vocabulary, algorithmic toolkit, and evaluation methodology required to
appreciate the experimental contributions of the subsequent chapters.  By
grounding our study in this foundation, we aim to make the design choices,
trade-offs, and results of our GRPO scaling experiments both transparent
and reproducible.

\end{document}
